%%%%%%%%%%%%%%%%%%%%%%%%%%%%%%%%%%%%%%%%%%%%%%%%%%%%%%%%%%%%%%%%%%%%%%%%%%%%%%%%%%%%%%%%%%%%%%%%%%%%%%
%
%   Filename    : abstract.tex 
%
%   Description : This file will contain your abstract.
%                 
%%%%%%%%%%%%%%%%%%%%%%%%%%%%%%%%%%%%%%%%%%%%%%%%%%%%%%%%%%%%%%%%%%%%%%%%%%%%%%%%%%%%%%%%%%%%%%%%%%%%%%

\begin{abstract}
%% \begin{mdframed} [style=highlight]
%% \end{mdframed}


Ang larangan ng Automatic Personality Recognition (APR) para sa mga Pilipino ay patuloy na lumalago, ngunit karamihan sa mga pag-aaral ay nakatuon lamang sa tekstuwal na pagsusuri mula sa mga plataporma tulad ng X (dating Twitter), kahit ang mga ebidensiya tungkol sa pagsasama ng biswal at lingguwistikong pahiwatig ay makabuluhang nagpapabuti sa katumpakan ng paghula ng personalidad. Layunin ng pag-aaral na ito ay tugunan ang kakulangan ng multimodal na pananaliksik sa kontekstong Pilipino sa pamamagitan ng pagbuo ng isang APR framework para sa mga gumagamit ng Instagram, na pinagsasama ang mga katangian mula sa larawan, bilingual (Ingles/Tagalog) na teksto, at metadata ng account. Gamit ang mga classification model tulad ng Support Vector Machine (SVM) at Extreme Gradient Boosting (XGBoost) na may feature-level fusion, inihahambing ng pag-aaral ang bisa ng unimodal at multimodal na paraan sa limang malalaking katangian ng personalidad (Big Five) gamit ang PagkataoKo dataset. Ipinapakita ng mga resulta na ang predictive performance ay lubos na nakadepende sa katangian: ang mga tekstuwal na feature tulad ng Term Frequency-Inverse Document Frequency (TF-IDF) at Word2Vec ang pangunahing nagtutulak para sa Conscientiousness, Agreeableness, at Neuroticism, samantalang ang mga biswal na feature tulad ng Residual Network (ResNet) at Hue, Saturation, Value (HSV) ang nangingibabaw para sa Openness. Para sa Extraversion, ang pinakamahusay na performance ay nakita sa paggamit ng Word2Vec Posts Weighted Equally (PWE) na may SVM model, at Word2Vec Words Weighted Equally (WWE) na may Logistic Regression (LR) model. Higit pa rito, bagama't sa pangkalahatan ay nagpapahusay ng robustness ang multimodal integration at ang metadata ay nagbibigay ng mahalagang pagpipino para sa mga regularidad ng pag-uugali, ang ilang partikular na katangian ay nananatiling pinakamabisang nahuhula sa pamamagitan ng tiyak na unimodal na mga channel.


%
%  Do not put citations or quotes in the abract.
%
\end{abstract}
